\documentclass[11pt]{article}
\usepackage[a4paper,margin=24mm]{geometry}
\usepackage{amsmath,amssymb,hyperref}
\title{Divergent viscosity paths in conjecture 00000008834}
\author{ziangni-sys}
\date{4 October 2026}
\begin{document}
\maketitle
The source claims that viscosity solution trajectories always converge
to the zero set of a monotone inclusion. Neither version assumes that
this zero set is nonempty. We exhibit actual divergent regularized
solutions and an actual divergent viscosity flow, rather than merely
pointing to an empty target set. The rate clause is not needed.

On the real Hilbert line let $A(x)=\{1\}$ for every $x$. This graph is
monotone, since any two of its points have the same second coordinate.
It is also maximal monotone: if a monotone graph extension contains
$(x,u)$, comparison with $(x-1,1)$ and $(x+1,1)$ gives respectively
$u\ge1$ and $u\le1$, hence $u=1$. Thus every monotone extension equals
the original graph. Clearly $A$ has no zero.

\paragraph{Actual Tikhonov solutions.}
For every $\varepsilon>0$ the regularized inclusion
\[
0\in A(x)+\varepsilon x
\]
is exactly $1+\varepsilon x=0$ and has the unique solution
$x_\varepsilon=-1/\varepsilon$. Choose the genuine positive vanishing
parameter
\[
\varepsilon(t)=\frac1{t+1},\qquad t\ge0.
\]
Then its solution trajectory is $v(t)=-(t+1)$, which tends to
$-\infty$ as $t\to\infty$. This parameter ranges over $(0,1]$, so it
also describes the full small-positive-parameter solution path.

\paragraph{Actual regularized flow.}
For the customary continuous viscosity dynamics
\[
u'(t)+A(u(t))+\varepsilon(t)u(t)\ni0
\]
there is an explicit solution
\[
u(t)=-\frac{t+1}{2},\qquad u(0)=-\frac12.
\]
Indeed $u'(t)=-1/2$ and
\[
u'(t)+1+\varepsilon(t)u(t)
=-\frac12+1-\frac12=0
\]
for every $t\ge0$. This smooth solution also tends to $-\infty$.
In particular neither actual trajectory converges to any finite real
number, in norm or weakly (the identity continuous linear functional
already detects the divergence). Both common meanings of the viscosity
trajectory therefore fail the unqualified convergence claim. The
missing zero-existence hypothesis is essential; the example makes no
claim against a theorem that includes it.

\paragraph{Lean correspondence.}
The project uses the actual real inner product and proves maximality
against every monotone graph extension. It proves the full regularized
inclusion and unique solution for every positive parameter, positivity
and convergence to zero of $\varepsilon(t)$, actual derivatives and the
regularized differential inclusion on $t\ge0$, and actual nonconvergence
of both real-time trajectories. Nonconvergence is verified via their
affine integer samples and uniqueness of limits, not assumed from the
absence of zeros. Lean 4.19.0 and Mathlib commit
\texttt{c44e0c8ee63ca166450922a373c7409c5d26b00b} are pinned.
\end{document}
